\documentclass[11pt]{article}
\usepackage[margin=1in]{geometry}
\usepackage{amsmath,amssymb,amsthm}
\usepackage[hidelinks]{hyperref}
\newtheorem{theorem}{Theorem}
\newtheorem{lemma}[theorem]{Lemma}
\newtheorem{corollary}[theorem]{Corollary}
\title{Two initial moments at arbitrary lower exponent and the polynomial boundary condition}
\author{Internal elementary proof and exact replay; no priority claim}
\date{17 November 2026}
\begin{document}
\maketitle
For $d\ge1$ an integer and $\gamma>0$, put
\[
 T_{n,\gamma,d}(t)=\sum_{k=0}^n\binom nk(\gamma)_{dk}t^k.
\]
The shape varies with the moment order when both exponents are nonzero.
This note proves a uniform two--three moment separation at every lower
exponent, and identifies the polynomial condition selecting the solution
of the higher-gap parameter recurrence. It does not prove all-index
binomial Strong Factorial.

\section{A parameter-monotone initial separation}
\begin{lemma}\label{lem:inequality}
For every integer $d\ge1$ and real $y\ge x>0$,
\[
 \frac{(y)_{3d}}{(x)_{2d}}>3(y)_d-2(x)_d.                 \tag{1}
\]
\end{lemma}
\begin{proof}
Write $r(s)=(s)_d$. Its polynomial coefficients are nonnegative,
$r'(s)>0$ and $r''(s)\ge0$ for $s>0$. In particular $r$ is strictly
increasing and $r'$ is nondecreasing there. Put
\[
 H(y)=r(y)r(y+d)r(y+2d)
       -3r(x)r(x+d)r(y)+2r(x)^2r(x+d).
\]
At $y=x$,
\[
 H(x)=r(x)r(x+d)\bigl(r(x+2d)-r(x)\bigr)>0.
\]
For $y\ge x$, the product rule and the monotonicity of $r'$ give
\[
 \begin{split}
 \frac{d}{dy}\bigl(r(y)r(y+d)r(y+2d)\bigr)
 \ge r'(y)\bigl(&r(y+d)r(y+2d)+r(y)r(y+2d)\\
               &+r(y)r(y+d)\bigr).
 \end{split}
\]
Each of the three products in parentheses is at least
$r(x)r(x+d)$, and the first two are strictly greater.
Thus $H'(y)>0$. Consequently $H(y)>0$, which is (1).
\end{proof}

\begin{theorem}\label{thm:two-three}
For every integer $d\ge1$ and all real $y\ge x>0$, the polynomials
$T_{2,x,d}$ and $T_{3,y,d}$ have no common complex zero.
\end{theorem}
\begin{proof}
The quadratic
\[
 T_{2,x,d}=1+2(x)_d t+(x)_{2d}t^2
\]
has negative discriminant, since
$(x)_{2d}=(x)_d(x+d)_d>(x)_d^2$. A common zero therefore
brings its distinct conjugate, and would imply
\[
 T_{3,y,d}=(1+ct)T_{2,x,d}
\]
for a real $c$. Comparing the leading coefficient gives
$c=(y)_{3d}/(x)_{2d}$, while comparing the linear coefficient gives
$c=3(y)_d-2(x)_d$. Lemma~\ref{lem:inequality} contradicts these
equalities.
\end{proof}

For the normalized positive Gamma functional
$\mathcal L_\beta(X^j)=(\beta)_j$, $\beta>0$, the moments of
$f=X^a(A+BX^d)$ satisfy, when $A\ne0$,
\[
 \mathcal L_\beta(f^n)
   =A^n(\beta)_{an}T_{n,\beta+an,d}(B/A).                   \tag{2}
\]
\begin{corollary}\label{cor:initial}
Let $a\ge0$ and $d\ge1$ be integers, and let $\beta>0$.
Every nonzero $X^a(A+BX^d)$ has no simultaneous zero Gamma moments
at orders two and three. It also has no simultaneous zero Gamma moments
at orders one and two. Thus every consecutive pair with terminal order
at most three contains a nonzero moment, for all lower exponents and gaps.
\end{corollary}
\begin{proof}
For orders two and three, apply Theorem~\ref{thm:two-three} with
$x=\beta+2a$ and $y=\beta+3a$. For orders one and two, the sole
root of $T_{1,\beta+a,d}$ is real, whereas
$T_{2,\beta+2a,d}$ has no real root by the discriminant calculation above.
The order-zero moment is one. If either $A$ or $B$ is zero, every
moment is a nonzero monomial moment.
\end{proof}
For ordinary factorials set $\beta=1$. No restriction on the complex
coefficients $A,B$ beyond nonzeroness of $f$ is imposed.

\section{The boundary condition at infinity for the parameter recurrence}
For terminal order $N\ge2$ define
\[
 U_N^{(d)}(s,z)=\sum_{k=0}^N(-1)^k\binom Nk(s+1)_{dk}z^{N-k},
 \qquad V(s)=U_{N-1}^{(d)}(s,z).
\]
The exact identities in the
\href{UNIFORM_PADE_PARABOLA_SFC.tex}{parameter-recurrence note}
give
\[
 z\bigl(V(s)-V(s+1)\bigr)
 =(s+2)_{d-1}\bigl((s+dN+1)V(s+d)
                    -(s+d+1)V(s+d+1)\bigr).              \tag{3}
\]
At a hypothetical common moment zero, setting
$l=\beta+a(N-1)-1$, $M=l+a$ gives the two finite conditions
\[
 V(l)=0,\qquad zV(M)=(M+1)_dV(M+d).                        \tag{4}
\]
These two conditions alone do not determine the order-$d+1$ difference
solution when $d\ge2$.

\begin{theorem}\label{thm:polynomial}
Fix $N\ge2$, $d\ge1$ and any complex $z$. The space of polynomial
solutions $V\in\mathbb C[s]$ of (3) is one-dimensional, spanned by
$U_{N-1}^{(d)}(s,z)$. Every nonzero such solution has degree
$m=d(N-1)$. The normalization
\[
 [s^m]V=(-1)^{N-1}                                       \tag{5}
\]
therefore selects $U_{N-1}^{(d)}$ uniquely.
Moreover this selected solution has the exact degree-gap condition
\[
 V(s)=(-1)^{N-1}(s+1)_m+Q(s),
 \qquad \deg Q\le m-d.                                   \tag{6}
\]
\end{theorem}
\begin{proof}
Let $V$ be nonzero of degree $j$ with leading coefficient $c$.
In the expression in parentheses on the right of (3), the coefficient
of $s^j$ is
\[
 c\bigl(d(N-1)-j\bigr).
\]
Indeed the degree-$j+1$ terms cancel, while the shifts in the
degree-$j$ terms contribute $-cj$.
If $j\ne d(N-1)$, the right side has degree $j+d-1$ and nonzero
leading coefficient. The left side has degree at most $j-1$.
This is impossible. Hence every nonzero polynomial solution has
degree $m=d(N-1)$.

The displayed coefficient formula for $U_{N-1}^{(d)}$ verifies that
it solves (3) and has degree $m$ and leading coefficient
$(-1)^{N-1}$. If two polynomial solutions were independent,
a linear combination would cancel their degree-$m$ terms and give
a nonzero polynomial solution of smaller degree, a contradiction.
This proves uniqueness and (5). Finally, the top summand $k=N-1$
in that coefficient formula is $(-1)^{N-1}(s+1)_m$, while all other
summands have degree at most $d(N-2)=m-d$, proving (6).
\end{proof}

Polynomiality supplies the missing selection condition for this
particular recurrence solution. It does not prove that (4) is
incompatible with it, and does not supply a positive finite-matrix or
zero-free-region argument. The all-index lower-exponent problem
remains open. The distinction is essential: recurrence order alone
does not license arbitrary initial data, while a uniquely selected
polynomial solution may still satisfy two specified finite conditions.

The accompanying exact replay checks the coefficient identities and
the polynomial-space rank at bounded indices, plus the strict
inequality at rational parameters. The universal evidence is the
elementary proof above; the bounded replay does not establish the
quantifiers. Independent mathematical review, independent replay,
formal verification and external review are not claimed.
\end{document}
